% ------------
% | Preamble |
% ------------
%\pdfsuppresswarningpagegroup=1
%\pdfminorversion=7

\usepackage[ngerman,british]{babel}
%\setlength{\marginparwidth}{1.5cm}
\usepackage[textsize=scriptsize]{todonotes}
\usepackage{lipsum}
\usepackage[figuresright]{rotating}


\usepackage{titlecaps}
\Addlcwords{a, an, and, as, at, but, by, for, if, in, nor, of, off, on, or, per, so, the, to, up, via, yet}

% --------------
% | Custom |
% --------------
\usepackage{xltabular}
\usepackage{multicol}

% BPMS
\newcommand{\arti}{design theory}
\newcommand{\faarti}{functional affordance}
\newcommand{\faartipl}{functional affordances}
\newcommand{\toarti}{technical object}
\newcommand{\toartipl}{technical objects}
\newcommand{\company}{WomOut}

%\newcolumntype{M}[1]{>{\centering\arraybackslash}m{#1}}
\newcommand\tabrotate[1]{\begin{turn}{90}\rlap{#1}\end{turn}}

% --------------
% | Font Stuff |
% --------------
% Alternative to Myriad Pro
\usepackage{tgheros} % qhv

% Alternative to Minion Pro
\usepackage{tgpagella} % qpl

\usepackage[utf8]{inputenc}
\usepackage[LGR,T1]{fontenc} 			
 
% Greek Font
%\usepackage{palatino}
\usepackage{pifont}						% Pi font with some symbols
\usepackage[right]{eurosym}				% Euro-Symbol
\usepackage{tablefootnote}
\usepackage{hyphenat}
\usepackage{adjustbox}
\usepackage{minibox}
\usepackage{framed}
\usepackage[misc,geometry]{ifsym}
\usepackage[
	activate={true,nocompatibility},
	final,
	babel=true,
	tracking=true,
	kerning=true,
	spacing=true,
	factor=1100,
	stretch=10,
	shrink=10]{microtype}
\microtypecontext{spacing=nonfrench}
\SetTracking{encoding={*}, shape=sc}{40}
\SetExtraKerning[unit=space]
{encoding={*}, family={bch}, series={*}, size={footnotesize,small,normalsize}}
{\textendash={400,400}, % en-dash, add more space around it
	"28={ ,150}, % left bracket, add space from right
	"29={150, }, % right bracket, add space from left
	\textquotedblleft={ ,150}, % left quotation mark, space from right
	\textquotedblright={150, }} % right quotation mark, space from left

\emergencystretch=1.5em

% activate={true,nocompatibility} - activate protrusion and expansion
% final - enable microtype; use "draft" to disable
% tracking=true, kerning=true, spacing=true - activate these techniques
% factor=1100 - add 10% to the protrusion amount (default is 1000)
% stretch=10, shrink=10 - reduce stretchability/shrinkability (default is 20/20)

% ----------------
% | Hyphenation  |
% ----------------
\usepackage[shortcuts]{extdash}
\usepackage{hyphenat}	% enable hyphenation
%\usepackage[none]{hyphenat} % disable hyphenation

\hyphenation{me-di-um Ar-nott Mend-ling Poly-vyanyy leg-isla-tive}

\usepackage{url}
\makeatletter
\g@addto@macro{\UrlBreaks}{\UrlOrds}
\makeatother

% ----------------------------------
% | KOMA: Adjust Headings |
% ----------------------------------
\setkomafont{part}{\Huge\bfseries}
%\setkomafont{chapter}{\Huge}
%\setkomafont{section}{\LARGE}
\setkomafont{subsection}{\large\emph}
%\setkomafont{subsubsection}{\large}
%\setkomafont{paragraph}{\large}
%\setkomafont{subparagraph}{\large}

% ----------------------------------
% | KOMA: Adjust Headers |
% ----------------------------------
\usepackage[pagestyleset=standard,
            markcase=used]
            {scrlayer-scrpage} 
% pagestyleset=standard matches the page styles of the standard LaTeX document classes, i.e.
% (1) frontmatter uses Roman numerals
% (2) no running heads for part and chapter pages
% (3) running heads inner-aligned, page numbers outer-aligned
% (4) automatic running heads are converted to capital letters

\setkomafont{pagehead}{\sffamily} % Set header font family to sans serif
\setkomafont{pagenumber}{\sffamily} % Set pagenumber font family to sans serif

%%%
% Remove formatting symbols before marks
%%%

% \renewcommand*{\chaptermarkformat}{}
% \renewcommand*{\sectionmarkformat}{}

\renewcommand*{\indexpagestyle}{headings}
\renewcommand*{\partpagestyle}{empty}
\renewcommand*{\chapterpagestyle}{empty}

\renewcommand*{\partformat}{\thepart\autodot}

% ------------------------------
% | Compact Paragraph Sections |
% ------------------------------
\makeatletter
\renewcommand{\paragraph}{%
	\@startsection{paragraph}{4}%
	{\z@}{3.25ex \@plus 1ex \@minus .2ex}{-1em}%
	{\normalfont\normalsize\bfseries}%
}
\makeatother

% ----------------
% | ToC & Tweaks |
% ----------------
% fix spacing in LoT, e.g., 19.18
 \makeatletter
 \renewcommand*\l@table{\@dottedtocline{1}{1em}{2.8em}}
 \renewcommand*\l@figure{\@dottedtocline{1}{1em}{2.8em}}
 \let\l@table\l@figure
 \makeatother

% Adjustments to the style of the tables of contents, list of figures, ...
% Chapter
\DeclareTOCStyleEntry[raggedentrytext=on]{tocline}{chapter} % fixes overfull hboxes in chapter titles in toc
% Section
\DeclareTOCStyleEntry[raggedentrytext=on]{tocline}{section} % fixes overfull hboxes in section titles in toc
\DeclareTOCStyleEntry[entryformat=\sffamily]{tocline}{section} % force sans serif on section entries in toc
\DeclareTOCStyleEntry[pagenumberformat=\sffamily]{tocline}{section} % force sans serif on page number of section entries in toc
% Subsection
\DeclareTOCStyleEntry[entryformat=\sffamily]{tocline}{subsection} % force sans serif on subsection entries in toc
\DeclareTOCStyleEntry[raggedentrytext=on]{tocline}{section} % force sans serif on subsection entries in toc
\DeclareTOCStyleEntry[pagenumberformat=\sffamily]{tocline}{subsection} % force sans serif on page number of subsection entries in toc
% Figure
\DeclareTOCStyleEntry[entryformat=\sffamily]{tocline}{figure} % force sans serif on section entries in toc
\DeclareTOCStyleEntry[pagenumberformat=\sffamily]{tocline}{figure} % force sans serif on page number of section entries in toc
\DeclareTOCStyleEntry[entryformat=\sffamily]{tocline}{table} % force sans serif on section entries in toc
\DeclareTOCStyleEntry[pagenumberformat=\sffamily]{tocline}{table} % force sans serif on page number of section entries in toc
%\addtokomafont{disposition}{\rmfamily}			% set chapter font, default in KOMA-Script is sffamily 
\addtokomafont{dictum}{\sffamily\itshape\justify}			% set dictum font
\addtokomafont{partentrypagenumber}{\nullfont}	% no page numbers for parts
\renewcommand*{\dictumwidth}{.8\textwidth}
\renewcommand*{\dictumrule}{\smallskip}
\renewcommand*{\raggeddictumtext}{}
\renewcommand*{\dictumauthorformat}[1]{#1\vspace{1em}}

% ----------------
% | Tight Enumerate |
% ----------------
\newenvironment{tight_enumerate}{
	\begin{enumerate}
		\setlength{\itemsep}{0pt}
		\setlength{\parskip}{0pt}
	}{\end{enumerate}}

% -------------
% | Tight Itemize |
% -------------
\newenvironment{tight_itemize}{
	\begin{itemize}
		\setlength{\parskip}{0pt}
	}{\end{itemize}}

% -------------
% | Float Env |
% -------------
%\usepackage{flafter}  			% figures and tables always after first reference

\usepackage[section]{placeins} 	% Floatbarrier
\usepackage[bottom]{footmisc}	% no figures below footnotes
%\usepackage{float}				% use H to force placement

% ---------------------------------------------
% | Do not reset footnote counter per chapter |
% ---------------------------------------------
\usepackage{chngcntr}
\counterwithout{footnote}{chapter}
\deffootnote{2em}{1em}{\textsuperscript{\thefootnotemark}\hspace{0.5em}}

\usepackage{graphicx}
\usepackage{epstopdf}
\usepackage{enumerate}
\usepackage{soul}
% ------------
% | Geometry |
% ------------
\usepackage[includehead,
			top 		= 2.5cm,
			bottom 		= 2.5cm,
			inner 		= 2.5cm,
			outer		= 2.5cm,
			%left 		= 2.8cm,		% old: 2.8cm
			%right 		= 2.8cm,		% old: 2.8cm
			footskip	= 1cm]{geometry}

% -----------------------
% | Theorem Environment |
% -----------------------
\usepackage{amsmath}
\usepackage{amssymb}
\usepackage{amsthm}
\theoremstyle{plain}
\newtheorem{obs}{Observation}[] % reset theorem numbering for each chapter
\newtheorem*{reo}{Research Objective} % * disables numbering of theorem. Remove if several research objectives.
\newtheorem{req}{Research Question}[] % reset theorem numbering for each chapter
\newtheorem{hypothesis}{Hypothesis}
\usepackage[euler]{textgreek}

% Definitions
\theoremstyle{definition}
\newtheorem{definition}{Definition}[]

\usepackage{letltxmacro}

% --------
% | tikz |
% --------
\usepackage{tikz}
\usetikzlibrary{calc}
\usetikzlibrary{matrix}
\usetikzlibrary{positioning}
\usetikzlibrary{shapes.geometric}
%
\usepackage{pgfplots}
\pgfplotsset{width=0.9\textwidth,compat=1.13}
\newcommand*{\escape}[1]{\texttt{\textbackslash#1}}
\usepackage{forest}
\usepackage{tikz-dependency}
\usepackage{lstautogobble}
% ------------
% | Acronyms |
% ------------
% Attention: works only with version 2 of acro package. Add [version=2] as option to usepackage if compiling fails.
\usepackage{acro}[=v2]
\acsetup{
    extra-style=comma,
    list-style=longtable, 
    list-short-format={\bfseries}
}
% Returning to version 2 requires the following code to be added to suppress an error
\acsetup{display-foreign = false}
\acsetup{foreign-format = {\scshape} }
\ExplSyntaxOn
\prop_new:N \l__acro_foreign_format_prop
\ExplSyntaxOff

\usepackage{expl3}
\ExplSyntaxOn
\tex_let:D \c_minus_one \scan_stop:
\int_const:Nn \c_minus_one {-1}
\ExplSyntaxOff

\usepackage{etoolbox}
\usepackage{paralist}	 	% compact itemize lists

\usepackage{makeidx}
\usepackage{ellipsis}		% Improve whitespace around \ldot


% ----------------
% | Line Spacing |
% ----------------
\usepackage{setspace}
\onehalfspacing

% ------------
% | ragged2e |
% ------------
\usepackage{ragged2e}		% Flattersatz mit Silbentrennung
\raggedbottom

% ----------
% | Tables |
% ----------
\usepackage{tabularx}
\usepackage{booktabs}
\usepackage{multirow}
\usepackage{colortbl}
\usepackage{tcolorbox}
\tcbuselibrary{breakable,skins}
\usepackage{longtable}
\usepackage{lscape}

% -----------
% | Caption |
% -----------

% The caption package offers customization of captions in floating environments such figure and table and cooperates with many other packages.

\usepackage[font=sf]{caption}
%\captionsetup[algorithm]
%             {format=hang,
%              labelfont={bf},
%              aboveskip=2ex,
%              belowskip=1ex,
%              justification=raggedright, % Horizontal zentriert
%              singlelinecheck = false,
%              position=bottom,           % Vertikal "drunter"
%              textformat=simple}         %

\captionsetup{format=hang,
        font={sf,small},
	labelfont={bf},
	aboveskip=2ex,
	belowskip=1ex,
	justification=raggedright, % Horizontal zentriert
	singlelinecheck = false,
	position=bottom,           % Vertikal "drunter"
	textformat=simple}         %

\usepackage{subcaption}
\newcommand{\source}[1]{
	\caption*{\footnotesize{\hfill \textsc{Source}: {#1}}} \vspace{-4mm} 
}
\newcommand{\quoteSource}[1]{
	\nopagebreak\vspace{-2mm}
	\footnotesize{\hfill \textsc{Source}: {#1}}
}

% --------------
% | Algorithms |
% --------------

\usepackage{algorithmicx}
\usepackage[plain]{algorithm}% http://ctan.org/pkg/algorithms
\usepackage[noend]{algpseudocode}% http://ctan.org/pkg/algorithmicx

\algrenewcommand\alglinenumber[1]{\tiny #1:}
\algnewcommand\algorithmicswitch{\textbf{switch}}
\algnewcommand\algorithmiccase{\textbf{case}}
\algnewcommand\algorithmicassert{\texttt{assert}}
\algnewcommand\Assert[1]{\State \algorithmicassert(#1)}%
\algnewcommand{\IfThen}[2]{% \IfThenElse{<if>}{<then>}{<else>}
	\State \algorithmicif\ #1\ \algorithmicthen\ #2}
% New "environments"
\algdef{SE}[SWITCH]{Switch}{EndSwitch}[1]{\algorithmicswitch\ #1\ \algorithmicdo}{\algorithmicend\ \algorithmicswitch}%
\algdef{SE}[CASE]{Case}{EndCase}[1]{\algorithmiccase\ #1}{\algorithmicend\ \algorithmiccase}%
\algtext*{EndSwitch}%
\algtext*{EndCase}%
\algnewcommand{\algorithmicgoto}{\textbf{go to}}%
\algnewcommand{\Goto}[1]{\algorithmicgoto~\ref{#1}}%


% -------------
% | ToC Depth |
% -------------
% -->  bis (2) subsection | (3) subsubsection | (4) paragraph
\setcounter{secnumdepth}{2}		% Numbering in documente
\setcounter{tocdepth}{2}		% in ToC

% ------------
% | Listings |
% ------------
\usepackage{listings}
	\lstset{linewidth=\linewidth}
	\lstset{breaklines=true}
	\lstset{breakatwhitespace=true}
	\lstset{language=XML}
	\lstset{usekeywordsintag=true}
	\lstset{basicstyle=\scriptsize\ttfamily}
	\lstset{stringstyle=\ttfamily}
	\lstset{keywordstyle=\sffamily}
	\lstset{numbers=left, numberstyle=\tiny, numbersep=5pt}


% ---------------------
% | Tab Item in Table |
% ---------------------
\newcommand{\tabitem}{~~\llap{\textbullet}~~}

% ----------------
% | Circled Mark |
% ----------------
\newcommand\circledmark[1][gray!50]{%
  \tikz[baseline=(A.south)]{
    \node[circle,fill=#1,inner sep=0.75ex] (A) {};
    \node at (A) {\ding{51}};
  }%
}

\newcommand{\xmark}{\ding{55}}%

\newcommand\circledcross[1][gray!50]{%
	\tikz[baseline=(A.south)]{
		\node[circle,fill=#1,inner sep=0.75ex] (A) {};
		\node at (A) {\ding{55}};
	}%
}

% ----------------------------------
% | Hurenkinder und Schusterjungen |
% ----------------------------------
\clubpenalty=10000
\widowpenalty=10000
\displaywidowpenalty=10000
\interfootnotelinepenalty=10000

% ----------------
% | Bibliography |
% ----------------
\usepackage[backend=biber,giveninits=true,doi=true,url=true,isbn=false,natbib=true,style=apa,maxbibnames=99,maxcitenames=2,hyperref=true,backref=false,defernumbers=true,dashed=false,autopunct=true,uniquename=false,uniquelist=false,useprefix=true,date=long,sorting=nyt]{biblatex}
\usepackage[autostyle=true,german=quotes]{csquotes}
\usepackage{breakcites}
\apptocmd{\sloppy}{\hbadness 10000\relax}{}{} %prevent underfull hbox warnings

\renewcommand{\finentrypunct}{.}


% -------------------------------------
% | Rename Bibliography to References |
% -------------------------------------
\DefineBibliographyStrings{english}{%
	bibliography = {References},
}

% ------------
% | Citeapos |
% ------------
\def\citeapos#1{\citeauthor{#1}'s (\citeyear{#1})}

% --------------------
% | Tweak Full-Cite  |
% --------------------
\preto\fullcite{\AtNextCite{\defcounter{maxnames}{99}}}

% -----------------
% | Table Columns |
% -----------------
\newcolumntype{L}[1]{>{\raggedright\arraybackslash}p{#1}}   % linksbündig mit Breitenangabe
\newcolumntype{C}[1]{>{\centering\arraybackslash}p{#1}}     % zentriert mit Breitenangabe
\newcolumntype{R}[1]{>{\raggedleft\arraybackslash}p{#1}}    % rechtsbündig mit Breitenangabe

\newcolumntype{Y}{>{\raggedright\arraybackslash}X}




\newcommand{\ks}[1]{\hangindent=0.5cm \textbf{#1}}

% ----------------------
% | Abstract in Part B |
% ----------------------
\newcommand\paperTitle[1]{%
	\newpage
	\begin{center}
	{\Large\sffamily\noindent
	\textbf{#1}}
	\end{center}
	\parskip
}

\newenvironment{abstract}{%
  \vskip-\lastskip
  \quotation
  \vskip-\partopsep\vskip-\parsep\vskip-\itemsep\vskip-2\baselineskip\vskip-\parskip%
  \leavevmode\small\subparagraph*{\abstractname.}%
}{%
  \par\endquotation
}


% -----------------------------------
% | Define mytexttt with hypenation |
% -----------------------------------
\makeatletter
\DeclareRobustCommand\ttfamily
        {\not@math@alphabet\ttfamily\mathtt
         \fontfamily\ttdefault\selectfont\hyphenchar\font=-1\relax}
\makeatother
\DeclareTextFontCommand{\mytexttt}{\ttfamily\hyphenchar\font=45\relax}



% ------------------
% | Custom colors  |
% ------------------
\usepackage{xcolor}
\definecolor{LightGray}{gray}{0.9}
\definecolor{mygreen}{rgb}{0,0.6,0}
\definecolor{mygray}{rgb}{0.5,0.5,0.5}
\definecolor{mymauve}{rgb}{0.58,0,0.82}
\definecolor{mygray}{gray}{0.85}
\definecolor{mycyan}{rgb}{0.88,1,1}
\definecolor{LightGray}{gray}{0.9}
\definecolor{crowdstrom}{RGB}{46, 170, 225}
\definecolor{mygreen}{rgb}{0,0.6,0}
\definecolor{mygray}{rgb}{0.5,0.5,0.5}
\definecolor{mymauve}{rgb}{0.58,0,0.82}
\definecolor{mygray}{gray}{0.85}
\definecolor{mycyan}{rgb}{0.88,1,1}
\colorlet{punct}{red!60!black}
\definecolor{background}{HTML}{EEEEEE}
\definecolor{delim}{RGB}{20,105,176}
\colorlet{numb}{magenta!60!black}
\definecolor{colorPositive}{RGB}{0,100,0}
\definecolor{colorNeutral}{RGB}{100,0,0}
\definecolor{colorNegative}{RGB}{95,95,95}

% ------------------
% | Abbreviations  |
% ------------------
% For correct spacing after abbreviations
\usepackage{xspace}
\DeclareRobustCommand\xdot{.\@\xspace}
\newcommand{\etc}{\mbox{etc}\xdot}
\newcommand{\eg}{\mbox{e.g}\xdot}
\newcommand{\Eg}{\mbox{E.g}\xdot}
\newcommand{\etal}{\mbox{et al}\xdot}
\newcommand{\esp}{\mbox{esp}\xdot}
\newcommand{\ie}{\mbox{i.e}\xdot}
\newcommand{\wrt}{\mbox{w.r.t}\xdot}
\newcommand{\sect}{\mbox{sec}\xdot}
\newcommand{\resp}{\mbox{resp}\xdot}
\newcommand{\apprx}{\mbox{approx}\xdot}
\newcommand{\tab}{\mbox{tab}\xdot}
\newcommand{\Tab}{\mbox{Tab}\xdot}
\newcommand{\fig}{\mbox{fig}\xdot}
\newcommand{\Fig}{\mbox{Fig}\xdot}
\newcommand{\viz}{\mbox{viz}\xdot}
\newcommand{\Ie}{\mbox{I.e}\xdot}
\newcommand{\cf}{\mbox{cf}\xdot}
\newcommand{\as}{\mbox{a.s}\xdot}
\newcommand{\iid}{\mbox{i.i.d}\xdot}
\newcommand{\p}[1]{\mbox{p.~#1}}
\newcommand{\pp}[1]{\mbox{pp.~#1}}
\newcommand{\pf}[1]{\mbox{p.~#1~f}\xdot}
\newcommand*\elide{\textup{[\,\dots]}}


\newenvironment{myquote}%
{\list{}{\leftmargin=0.3in\rightmargin=0.3in}\item[]}%
{\endlist}

% ------------
% | Hyperref |
% ------------
\PassOptionsToPackage{hyphens}{url}
\usepackage[breaklinks=true,
			plainpages=false,
			linktocpage=true,
			hidelinks,
			unicode]{hyperref}
\hypersetup{
	pdfauthor={\myname},
	pdftitle={\myplaintitle},
	pdfsubject={Dissertation},
	pdflang=EN}

	% -------------
	% | CleverRef |
	% -------------
	\usepackage{cleveref} % use same footnote mulitple times
	\crefformat{footnote}{#2\footnotemark[#1]#3}

% --------------------------------
% | scrhack: must be loaded last |
% --------------------------------
\usepackage{scrhack}

% % To create list of own publications by year
% %Pubs 

\usepackage{xparse}
% Taken from https://tex.stackexchange.com/questions/472588/sort-all-citations-from-one-year-automatically-under-a-heading-with-the-year-num
\ExplSyntaxOn
\seq_new:N \g__blxbibbyyear_yearlist_seq

\cs_new:Npn \__blxbibbyyear_seq_gput_right_once:Nn #1 #2
{
  \seq_if_in:NnF #1 {#2}
  { \seq_gput_right:Nn #1 {#2} }
}

\cs_generate_variant:Nn \__blxbibbyyear_seq_gput_right_once:Nn { NV, Nx }


\prg_new_conditional:Nnn \blx_field_if_undef:n { p, T, F , TF }
{
  \use:c { iffieldundef } { #1 } { \prg_return_true: } { \prg_return_false: }
}

% unfortunately, \iffieldint is not expandable, so no p version for us, boo
\prg_new_protected_conditional:Nnn \blx_field_if_int:n { T, F , TF }
{
  \iffieldint { #1 } { \prg_return_true: } { \prg_return_false: }
}

\AtDataInput
{
  \blx_field_if_undef:nF { labeldatesource }
  {
    \blx_field_if_undef:nTF { \thefield{labeldatesource}year }
    {
      \blx_field_if_undef:nF { \thefield{labeldatesource} }
      {
        \blx_field_if_int:nT { \thefield{labeldatesource} }
        {
          \__blxbibbyyear_seq_gput_right_once:Nx \g__blxbibbyyear_yearlist_seq
          { \thefield{\thefield{labeldatesource}} }
        }
      }
    }
    {
      \__blxbibbyyear_seq_gput_right_once:Nx \g__blxbibbyyear_yearlist_seq
      { \thefield{\thefield{labeldatesource}year} }
    }
  }
}

\cs_new:Npn \blxbibbyyear_seq_sort_bydirection:NN #1 #2
{
  \seq_sort:Nn #2
  {
    \int_compare:nNnTF { ##1 } #1 { ##2 }
    { \sort_return_swapped: }
    { \sort_return_same: }
  }
}

\cs_new:Nn \blxbibbyyear_seq_sort_descending:N
{
  \blxbibbyyear_seq_sort_bydirection:NN < #1
}

\cs_new_nopar:Npn \blxbibbyyear_print_yearbib:nn #1 #2
{

  \defbibcheck{thisyear}
  {
    \blx_field_if_int:nTF { labelyear }
    {
      \int_compare:nNnF { \thefield{labelyear} } = { #1 }
      { \skipentry }
    }
    { \skipentry }
  }
  % aktuell "hack" mit filter nach Jahr, da sonst leere Jahreseinträge angezeigt werden
  \ifnumless{#1}{\yearoffirstpub}{}{
  \printbibliography[heading=subbibliography, title=#1, check=thisyear, #2]
    }
}

\DeclareDocumentCommand{\printbibbyyear}{O{}}
{
  \blxbibbyyear_seq_sort_descending:N \g__blxbibbyyear_yearlist_seq
  \seq_map_inline:Nn \g__blxbibbyyear_yearlist_seq
  { \blxbibbyyear_print_yearbib:nn {##1} {#1} }
}
\ExplSyntaxOff

% Audrey’s approach to reverse numbering
% at https://tex.stackexchange.com/a/22770/105447

% Count total number of entries in each refsection
\AtDataInput{%
  \csnumgdef{entrycount:\therefsection}{%
    \csuse{entrycount:\therefsection}+1}}

% Print the labelnumber as the total number of entries in the
% current refsection, minus the actual labelnumber, plus one
\DeclareFieldFormat{labelnumber}{\mkbibdesc{#1}}
\newrobustcmd*{\mkbibdesc}[1]{%
  \number\numexpr\csuse{entrycount:\therefsection}+1-#1\relax}
